\documentclass[aspectratio=43]{beamer}
\usepackage{myPresentation}  
\usepackage{metalogo}

\title[Demonstration of \texttt{myPresentation}]{The \texttt{myPresentation} Beamer Theme}
\subtitle{A Comprehensive Demonstration and Guide to Beamer}
\author{Zhenran Guo\inst{1} \and DeepSeek\inst{2}}
\institute{
  \inst{1} School of Economics, Shanghai University of Finance and Economics \\
  \inst{2} DeepSeek AI
}
\date{\today}

\begin{document}

\begin{frame}
\titlepage
\end{frame}

\begin{frame}{Outline}
\tableofcontents
\end{frame}

\section{Core Features of the Theme}

\begin{frame}{Fonts and Typography}
This theme uses:
\begin{itemize}
    \item \textbf{English font}: \emphblue{Optima} (a humanist sans-serif typeface) for text.
    \item \textbf{Math font}: \emphblue{ArsenalMath-Sans} (a clean sans-serif math font) via \texttt{unicode-math}.
    \item \textbf{Color palette}: harmonious blue, red, and green for structure, alerts, and examples.
\end{itemize}
\pause
Example math in sans-serif style:
\[
\E[X] = \int_{\R} x\, dF(x), \quad \Var(Y) = \E[Y^2] - (\E[Y])^2.
\]
\end{frame}

\begin{frame}{Block Environments}
Three types of blocks are predefined:
\begin{block}{Standard Block}
This is a regular block with light blue background.
\end{block}
\pause
\begin{alertblock}{Alert Block}
Use it to highlight critical points.
\end{alertblock}
\pause
\begin{exampleblock}{Example Block}
Suitable for examples or positive remarks.
\end{exampleblock}
\end{frame}

\begin{frame}{Custom Box Environments}
The package also provides three handy environments with smaller fonts:
\begin{infobox}{Info Box}
A compact block, ideal for additional notes.
\end{infobox}
\pause
\begin{alertbox}{Alert Box}
Same as alertblock but with smaller text.
\end{alertbox}
\pause
\begin{examplebox}{Example Box}
For concise examples.
\end{examplebox}
\end{frame}

\begin{frame}{Custom Math Commands}
The theme provides convenient shortcuts:
\[
\R,\ \N,\ \C,\ \E,\ \Var,\ \Cov,\ \mat{A},\ \vct{x}.
\]
Example: $\mat{A}\vct{x} = \lambda \vct{x}$.

Also available: $\Prob$, $\tp$ (transpose), $\Z$, $\Q$, etc.

\note{This is a speaker's note. Use dual-screen mode to view notes.}
\end{frame}

\begin{frame}[allowframebreaks]{Using \texttt{allowframebreaks} for Long Content}
    This frame demonstrates the \texttt{allowframebreaks} option. When a frame has too much content, Beamer automatically splits it across multiple pages. The footline will show the frame number with a continuation suffix (e.g., 1/2, 2/2).

    \begin{exampleblock}{Key Points}
        \begin{itemize}
            \item Add \texttt{[allowframebreaks]} to the frame options.
            \item Use \texttt{longframe} environment as a shorthand (provided by the theme).
            \item Ideal for lengthy text, multiple items, or long lists.
        \end{itemize}
    \end{exampleblock}

    \begin{theorem}[Continuation Example]
        This theorem might be long and split across pages.
    \end{theorem}
    \begin{proof}
        Here is a very long proof that could span multiple pages if needed. \pause
        More content... \pause
        Even more... \pause
        And more... \pause
        Finally, we reach the end. \hfill $\square$
    \end{proof}

    \begin{itemize}
        \item Item 1
        \item Item 2
        \item Item 3
        \item Item 4
        \item Item 5
        \item Item 6
        \item Item 7
        \item Item 8
        \item Item 9
        \item Item 10
        \item Item 11
        \item Item 12
        \item Item 13
        \item Item 14
        \item Item 15
        \item Item 16
        \item Item 17
        \item Item 18
        \item Item 19
        \item Item 20
    \end{itemize}
\end{frame}

\section{Overlay Specifications in Beamer}

\begin{frame}{Basic Overlays: \texttt{\textbackslash pause}}
This text appears first.
\pause
This appears second.
\pause
This appears third.
\end{frame}

\begin{frame}{Item by Item with \texttt{\textbackslash pause}}
\begin{itemize}
    \item First item \pause
    \item Second item \pause
    \item Third item
\end{itemize}
\end{frame}

\begin{frame}{Using \texttt{\textbackslash only} and \texttt{\textbackslash onslide}}
\only<1>{This appears only on slide 1.}
\only<2>{This appears only on slide 2.}
\onslide<3>{This appears on slide 3 and stays.}
\end{frame}

\begin{frame}{More Complex Overlays with \texttt{< >}}
\begin{itemize}[<+-| alert@+>]
    \item This item will be highlighted when it appears.
    \item So will this one.
    \item And this one.
\end{itemize}
\end{frame}

\begin{frame}{Combining Overlays with Graphics}
\centering
\includegraphics<1>[width=0.3\textwidth]{example-image-a}%
\includegraphics<2>[width=0.3\textwidth]{example-image-b}%
\includegraphics<3>[width=0.3\textwidth]{example-image-c}

\onslide<1>{This is image A.}
\onslide<2>{This is image B.}
\onslide<3>{This is image C.}
\end{frame}

\begin{frame}{Text Replacement with \texttt{\textbackslash only}}
\only<1>{Version 1 of the text.}
\only<2>{Version 2 of the text.}
\only<3>{Version 3 of the text.}
\end{frame}

\begin{frame}{Incremental Uncovering with \texttt{\textbackslash uncover}}
\uncover<1->{This is visible from slide 1 onward.}
\uncover<2->{This is visible from slide 2 onward.}
\uncover<3->{This is visible from slide 3 onward.}
\end{frame}

\begin{frame}{Using \texttt{\textbackslash visible} and \texttt{\textbackslash invisible}}
\visible<1>{Only visible on slide 1.}
\visible<2>{Only visible on slide 2.}
\invisible<1>{This is invisible on slide 1 but appears on slide 2.}
\end{frame}

\begin{frame}{Overlay Specifications on Commands}
\textbf<1>{Bold on slide 1 only.}
\textbf<2>{Bold on slide 2 only.}
\textbf<3>{Bold on slide 3 only.}
\end{frame}

\section{Special Pages}

\begin{frame}{Footline and Header}
The footline displays the short title, current section, and frame number. The header shows navigation (if enabled). 
\end{frame}

\makethankyoupage[takeaway={Hope \texttt{myPresentation} Beamer Theme and the guide to layout help you with making professional slides!}]

\begin{frame}{Compilation Instructions}
    \begin{exampleblock}{To compile this presentation}
        \begin{itemize}
            \item Use \LuaLaTeX (required for \texttt{fontspec} and \texttt{unicode-math}).
            \item Run \LuaLaTeX \texttt{examplePre.tex} at least \textbf{twice} to resolve the table of contents and cross-references.
            \item For notes, use a PDF viewer with dual-screen support (e.g., \texttt{pympress}).
            \item For long frames, use \texttt{[allowframebreaks]} or the \texttt{longframe} environment.
        \end{itemize}
    \end{exampleblock}
\end{frame}

\iffalse
% Alternative thank-you page with contact info (uncomment to use)
\makethankyoupage[email =\href{mailto:john.doe@example.com}{john.doe@example.com},
  homepage =\href{https://johndoe.com}{johndoe.com},
  source =\href{https://github.com/johndoe}{github.com/johndoe},
  takeaway ={Beamer overlays make presentations dynamic and engaging!}
]
\fi

\end{document}